In der Zeit $t = \tkkO$ liefert das Kraftwerk die Energie
\begin{align}
 E &= \EkkF \\
   &= \Pkk\cdot\tkk \\
   &= \Ekk \approx \result{\EkkP{12}{2}}
\end{align}
Rechnet man Leistung mal Zeit direkt in Megawatt und Stunden, erhält man die Energie in Megawattstunden bzw. Gigawattstunden:
\begin{align}
 E &= \PkkO\cdot\tkkO = \Ekw \approx \result{\EkwP{0}{2}}
\end{align}
